\chapter{Аналитическая часть}

\section{Постановка задачи}
В лабораторной работе №2 необходимо провести сравнительный анализ рекурсивного и нерекурсивного алголритма. По варианту необходимо найти максимальное число в последовательности чисел.\newline 

\section{Рекурсия}
\textbf{Рекурсия} — это такой способ организации вычислений, при котором процедура или функция в ходе выполнения обращается сама к себе. Другими словами, рекурсия является методом определения функций (процедур), при котором определяемая функция применена в теле своего же собственного определения [1].
\textbf{Рекурсивный алгоритм} — алгоритм, определяемый через себя [2].

\section{Механизм выполнения рекурсивных вызовов}
Каждый вызов рекурсии 

\clearpage
